%% ----------------------------------------------------------------------------
%%  §4.7  Alcuni teoremi sulle funzioni continue
%% ----------------------------------------------------------------------------
\section{Alcuni teoremi sulle funzioni continue}
\label{sec:teoremi-continue}

Il teorema seguente è analogo al teorema della permanenza del segno per le successioni
(\cref{teo:permanenza-segno}) e alla sua versione per i limiti di funzioni
(\cref{teo:confronto-funzioni}(i)); la differenza è che ora la conclusione vale anche
\emph{nel} punto $x_{0}$.

\begin{teorema}[Permanenza del segno]\label{teo:permanenza-segno-continue}\index{teorema!della permanenza del segno}
Sia $f$ una funzione definita in un intervallo $A$ e continua in $x_{0}\in A$. Se
$f(x_{0})>0$, esiste un numero $\delta>0$ con la proprietà che $f(x)>0$ per ogni
$x\in A\cap(x_{0}-\delta,x_{0}+\delta)$.
\end{teorema}

\begin{schema}
\Ipotesi{$f$ continua in $x_{0}$; \ $f(x_{0})>0$.}
\Tesi{$\esiste\delta>0$: \ $f(x)>0$ \ $\perogni x\in A$, $\abs{x-x_{0}}<\delta$.}
\end{schema}

\begin{strategia}
Si sceglie una tolleranza che non permetta a $f(x)$ di raggiungere lo zero: con
$\epsilon=f(x_{0})/2$, la caratterizzazione (c) della continuità tiene $f(x)$ sopra
$f(x_{0})/2>0$ in tutto un intorno di $x_{0}$, punto $x_{0}$ compreso.
\end{strategia}

\begin{dimostrazione}
La dimostrazione si fa come per le successioni: dato che $f(x_{0})>0$, possiamo scegliere
$\epsilon=f(x_{0})/2$. Per la \cref{prop:caratt-continuita}(c) esiste quindi un numero
$\delta>0$ per cui $\abs{f(x)-f(x_{0})}<f(x_{0})/2$ per ogni $x\in A$ con
$\abs{x-x_{0}}<\delta$. Ciò equivale a $-f(x_{0})/2<f(x)-f(x_{0})<f(x_{0})/2$; in
particolare
\begin{equation}\label{eq:permanenza-continue}
  f(x)>f(x_{0})-\frac{f(x_{0})}2=\frac{f(x_{0})}2>0 . \qedhere
\end{equation}
\end{dimostrazione}

Molto importanti sono i tre teoremi che seguono: il teorema dell'esistenza degli zeri,
il teorema dell'esistenza dei valori intermedi e il teorema di Weierstrass sull'esistenza
del massimo e del minimo. Sono, nello schema della \cref{fig:schema-teoremi}, i primi
teoremi di esistenza che riguardano le funzioni.

\subsection{Il teorema dell'esistenza degli zeri}

\begin{teorema}[Esistenza degli zeri]\label{teo:zeri}\index{teorema!dell'esistenza degli zeri}
Sia $f$ una funzione continua in un intervallo $[a,b]$. Se $f(a)<0$ e $f(b)>0$, allora
esiste $x_{0}\in(a,b)$ tale che $f(x_{0})=0$.
\end{teorema}

\begin{schema}
\Ipotesi{$f$ continua in $[a,b]$; \ $f(a)<0<f(b)$.}
\Tesi{$\esiste x_{0}\in(a,b)$: \ $f(x_{0})=0$.}
\end{schema}

Naturalmente la tesi vale anche se $f(a)>0$ e $f(b)<0$ (basta applicare il teorema a
$-f$): il teorema dell'esistenza degli zeri vale supponendo che i valori $f(a)$, $f(b)$
siano di segno discorde. L'enunciato garantisce l'esistenza di \emph{almeno} uno zero,
non la sua unicità: gli zeri possono essere molti (\cref{fig:zeri}). La dimostrazione
del teorema, con il metodo di bisezione, è riportata nel \cref{sec:bisezione}.

\begin{figure}[htb]
\centering
\begin{tikzpicture}
\begin{axis}[width=0.92\textwidth, height=5cm, xmin=0, xmax=6.4, ymin=-1.9, ymax=2.6,
  axis lines=middle, xtick={0.5,6}, xticklabels={$a$,$b$}, ytick=\empty,
  xlabel={$x$}, ylabel={$y$},
  every axis x label/.style={at={(axis description cs:1,0.46)},anchor=west}]
  \addplot[curva, domain=0.5:6, samples=150] {0.15*(x-1.4)*(x-3.1)*(x-4.9)};
  \addplot[only marks, mark=*, mark size=1.6pt, blu] coordinates {(0.5,-1.5444) (6,2.2011)};
  \addplot[only marks, mark=*, mark size=1.6pt, rame] coordinates {(1.4,0) (3.1,0) (4.9,0)};
  \addplot[tratteggio] coordinates {(0.5,0) (0.5,-1.5444)};
  \addplot[tratteggio] coordinates {(6,0) (6,2.2011)};
  \node[font=\scriptsize, right] at (axis cs:0.55,-1.5444) {$f(a)<0$};
  \node[font=\scriptsize, left] at (axis cs:5.95,2.2011) {$f(b)>0$};
\end{axis}
\end{tikzpicture}
\caption{Una funzione continua in $[a,b]$ che assume agli estremi valori di segno
discorde attraversa l'asse $x$ almeno una volta: qui gli zeri sono tre. Il teorema
garantisce che ce n'è almeno uno, non dice quanti sono né dove si trovano.}
\label{fig:zeri}
\end{figure}

Tutte le ipotesi sono essenziali. Senza la continuità la tesi è falsa: la funzione che
vale $-1$ in $[-1,0)$ e $1$ in $[0,1]$ cambia segno senza annullarsi. È essenziale anche
che il dominio sia un intervallo: $\frac1x$ assume valori di segno discorde in $-1$ e in
$1$, è continua nel suo dominio $\R\setminus\{0\}$, ma non si annulla mai. Infine, come
vedremo nel \cref{sec:bisezione}, è essenziale l'assioma di completezza.

Per mostrare la portata del teorema, consideriamo come esempio le due equazioni nella
incognita $x$
\begin{align}
  &x^{3}+x-1=0, \label{eq:cubica}\\
  &e^{x}+x=0, \label{eq:trascendente}
\end{align}
che non rientrano tra le equazioni algebriche di primo e secondo grado di cui è facile
ricordare la formula risolutiva. La prima delle due equazioni è un'\emph{equazione
algebrica di terzo grado}, mentre la seconda è un'\emph{equazione trascendente}.
Procediamo per tentativi, assegnando a $x$ alcuni valori:
\begin{center}
\renewcommand{\arraystretch}{1.4}
\begin{tabular}{c|ccccc}
  $x$ & $-2$ & $-1$ & $0$ & $1$ & $2$\\ \hline
  $x^{3}+x-1$ & $-11$ & $-3$ & $-1$ & $1$ & $9$\\
  $e^{x}+x$ & $e^{-2}-2$ & $e^{-1}-1$ & $1$ & $e+1$ & $e^{2}+2$
\end{tabular}
\end{center}
Nel caso $f(x)=x^{3}+x-1$, abbiamo $f(0)<0$, $f(1)>0$. Poiché $f$ è un polinomio, e
quindi è continua, in base al teorema dell'esistenza degli zeri esiste un numero $x_{0}$
nell'intervallo $(0,1)$ tale che $f(x_{0})=0$; cioè $x_{0}$ è una soluzione
dell'equazione \eqref{eq:cubica}. Nel \cref{sec:bisezione} vedremo come calcolare
numericamente tale radice, e troveremo che $x_{0}=0{,}682327\ldots$ Nel secondo caso,
$f(x)=e^{x}+x$ è continua (\cref{teo:continuita-esponenziale}) e risulta
$f(-1)=\frac1e-1<0$, $f(0)=1>0$. Quindi esiste nell'intervallo $(-1,0)$ una radice
$x_{0}$ dell'equazione \eqref{eq:trascendente}; troveremo che $x_{0}=-0{,}567143\ldots$
In entrambi i casi la radice è unica, perché le due funzioni sono strettamente crescenti
(somme di funzioni strettamente crescenti): due zeri distinti $x_{1}<x_{2}$ darebbero
$0=f(x_{1})<f(x_{2})=0$.

Esiste una formula risolutiva per le equazioni di terzo grado (si ricordi la nota storica
su Cardano e Bombelli del \cref{sec:complessi}), che dà come soluzione reale
dell'equazione \eqref{eq:cubica} il numero
\begin{equation}\label{eq:cardano}
  x_{0}=\sqrt[3]{\frac12+\sqrt{\frac{31}{108}}}-\sqrt[3]{\sqrt{\frac{31}{108}}-\frac12}
  =0{,}682327\ldots
\end{equation}
Viceversa, non è nota alcuna formula risolutiva, in termini di funzioni elementari, per
l'equazione \eqref{eq:trascendente}.

\begin{notastorica}[Le equazioni che non si risolvono con una formula]
Dopo le formule per le equazioni di terzo e quarto grado, trovate nel Cinquecento da
Scipione del Ferro, Niccolò Tartaglia, Gerolamo Cardano e Ludovico Ferrari, i matematici
cercarono per quasi tre secoli una formula «per radicali» per l'equazione generale di
quinto grado. Paolo Ruffini (1799) e, in forma completa, Niels Henrik Abel (1824)
dimostrarono che una tale formula \emph{non esiste}; Évariste Galois chiarì poco dopo
quali equazioni particolari siano risolubili per radicali. Per le equazioni
trascendenti come la \eqref{eq:trascendente} la situazione è ancora più netta: la sua
soluzione si esprime con la \emph{funzione $W$ di Lambert}, $x_{0}=-W(1)$, ma $W$ è
definita proprio come inversa di una funzione, $w\mapsto we^{w}$, e non è una funzione
elementare. Il numero $W(1)=0{,}567143\ldots$ è detto \emph{costante omega}. Il teorema
degli zeri, insieme al metodo di bisezione, cambia la prospettiva: anche quando una
formula non esiste, l'esistenza della soluzione è garantita, e la soluzione si può
calcolare con la precisione che si desidera.
\end{notastorica}

Applichiamo ancora una volta il teorema dell'esistenza degli zeri per ritrovare, per
un'altra via, un risultato del \cref{cap:complementi}: l'esistenza della radice $n$-esima
(\cref{teo:radice-n}). \emph{Per ogni
$y_{0}\ge0$ esiste un numero reale $x_{0}\ge0$ soluzione dell'equazione}
\begin{equation}\label{eq:radice-zeri}
  x^{n}=y_{0} .
\end{equation}
Ricordiamo che, essendo la funzione $x^{n}$ strettamente crescente per $x\ge0$
(\cref{prop:potenza-crescente}), un tale numero $x_{0}$ è unico. Dimostriamo che
l'equazione ha soluzione: se $n=1$ basta $x_{0}=y_{0}$, e quindi possiamo supporre
$n\ge2$; se $y_{0}=0$ naturalmente è $x_{0}=0$, e se $y_{0}=1$ basta prendere
$x_{0}=1$. Negli altri casi ($y_{0}>0$, $y_{0}\neq1$) poniamo $f(x)=x^{n}-y_{0}$, che è
continua; risulta
$f(0)=-y_{0}<0$, e rimane da trovare un punto dove $f$ è positiva. Se $y_{0}<1$, allora
$f(1)=1-y_{0}>0$, e quindi esiste una radice $x_{0}$ nell'intervallo $(0,1)$. Se invece
$y_{0}>1$, allora $f(y_{0})=y_{0}^{n}-y_{0}=y_{0}\bigl(y_{0}^{n-1}-1\bigr)>0$, e quindi in
questo caso esiste una radice nell'intervallo $(0,y_{0})$. Nel \cref{cap:complementi}
avevamo ottenuto lo stesso risultato direttamente con l'estremo superiore; il teorema
degli zeri, che a sua volta poggia sull'assioma di completezza, incorpora una volta per
tutte quel ragionamento.

\subsection{Il teorema dei valori intermedi}

\begin{teorema}[Primo teorema dei valori intermedi]\label{teo:valori-intermedi-1}\index{teorema!dei valori intermedi}
Una funzione continua in un intervallo $[a,b]$ assume tutti i valori compresi tra $f(a)$
e $f(b)$.
\end{teorema}

\begin{schema}
\Ipotesi{$f$ continua in $[a,b]$; \ $y_{0}$ compreso tra $f(a)$ e $f(b)$.}
\Tesi{$\esiste x_{0}\in[a,b]$: \ $f(x_{0})=y_{0}$.}
\end{schema}

\begin{strategia}
Si «abbassa» il grafico di $y_{0}$: la funzione $g(x)=f(x)-y_{0}$ è continua e cambia
segno tra $a$ e $b$, e uno zero di $g$ è esattamente un punto in cui $f$ vale $y_{0}$
(\cref{fig:valori-intermedi}). Il teorema dei valori intermedi è dunque il teorema degli
zeri applicato a una funzione traslata.
\end{strategia}

\begin{dimostrazione}
Per semplificare le notazioni consideriamo il caso in cui $f(a)\le f(b)$ (l'altro caso
si ottiene applicando questo a $-f$). La tesi consiste nel provare che, qualunque sia
$y_{0}\in[f(a),f(b)]$, esiste $x_{0}\in[a,b]$ tale che $f(x_{0})=y_{0}$. Se $y_{0}=f(a)$ si
può porre $x_{0}=a$; analogamente, se $y_{0}=f(b)$, basta prendere $x_{0}=b$. Per trattare
il caso $y_{0}\in(f(a),f(b))$ consideriamo la funzione
\begin{equation}\label{eq:vi-g}
  g(x)=f(x)-y_{0} \qquad\perogni x\in[a,b],
\end{equation}
che è continua per il \cref{teo:operazioni-continue}; essendo $f(a)<y_{0}<f(b)$, risulta
\begin{equation}\label{eq:vi-segni}
  g(a)=f(a)-y_{0}<0, \qquad g(b)=f(b)-y_{0}>0 .
\end{equation}
Per il teorema dell'esistenza degli zeri esiste un numero $x_{0}\in(a,b)$ tale che
$g(x_{0})=0$, cioè $f(x_{0})=y_{0}$.
\end{dimostrazione}

\begin{figure}[htb]
\centering
\begin{tikzpicture}
\begin{axis}[width=0.85\textwidth, height=5cm, xmin=0, xmax=5.4, ymin=0, ymax=3.3,
  axis lines=left, xtick={0.6,5}, xticklabels={$a$,$b$}, ytick={0.7,1.552,2.404},
  yticklabels={$f(a)$,$y_{0}$,$f(b)$}, xlabel={$x$}, ylabel={$y$}]
  \addplot[curva, domain=0.6:5, samples=120] {0.7+0.32*(x-0.6)+0.95*sin(deg(1.5*(x-0.6)))};
  \addplot[curva2, domain=0:5.4] {1.552};
  \addplot[tratteggio] coordinates {(0,0.7) (0.6,0.7) (0.6,0)};
  \addplot[tratteggio] coordinates {(0,2.404) (5,2.404) (5,0)};
  \addplot[only marks, mark=*, mark size=1.6pt, rame] coordinates {(1.1332,1.552) (2.5277,1.552) (4.5017,1.552)};
\end{axis}
\end{tikzpicture}
\caption{Il teorema dei valori intermedi: ogni retta orizzontale $y=y_{0}$ con $y_{0}$
compreso tra $f(a)$ e $f(b)$ incontra il grafico di una funzione continua su $[a,b]$
almeno una volta (qui tre volte).}
\label{fig:valori-intermedi}
\end{figure}

\begin{idea}[L'equilibrio portanza--peso]
In volo rettilineo orizzontale uniforme la portanza deve uguagliare il peso:
$\frac12\rho V^{2}S\,C_{L}(\alpha)=W$. Per una data velocità $V$ e quota (cioè densità
$\rho$), esiste un angolo d'attacco che realizza l'equilibrio? Il coefficiente di
portanza $C_{L}(\alpha)$ è una funzione continua dell'incidenza (tra l'incidenza di
portanza nulla e lo stallo), e così la portanza $L(\alpha)=\frac12\rho V^{2}S\,C_{L}(\alpha)$.
Se a un'incidenza $\alpha_{1}$ la portanza è minore del peso e a un'incidenza $\alpha_{2}$
è maggiore, il teorema dei valori intermedi garantisce un'incidenza $\bar\alpha$ compresa
tra le due per cui $L(\bar\alpha)=W$. Se invece anche al $C_{L,\max}$ la portanza resta
inferiore al peso, l'equilibrio è impossibile: la velocità è minore della velocità di
stallo $V_{S}=\sqrt{2W/(\rho S\,C_{L,\max})}$.
\end{idea}

\subsection{Il teorema di Weierstrass}

\begin{teorema}[di Weierstrass]\label{teo:weierstrass}\index{teorema!di Weierstrass}
Sia $f$ una funzione continua in un intervallo chiuso e limitato $[a,b]$. Allora $f$
assume massimo e minimo in $[a,b]$, cioè esistono in $[a,b]$ due punti $x_{1}$, $x_{2}$
tali che
\begin{equation}\label{eq:weierstrass}
  f(x_{1})\le f(x)\le f(x_{2}) \qquad\perogni x\in[a,b].
\end{equation}
\end{teorema}

\begin{schema}
\Ipotesi{$f$ continua in $[a,b]$, intervallo chiuso e limitato.}
\Tesi{$\esiste x_{1},x_{2}\in[a,b]$: \ $f(x_{1})\le f(x)\le f(x_{2})$ \ $\perogni x\in[a,b]$.}
\end{schema}

I numeri $x_{1}$, $x_{2}$ sono detti rispettivamente \emph{punti di minimo} e \emph{di
massimo} per $f$ nell'intervallo $[a,b]$; i corrispondenti valori $m=f(x_{1})$ e
$M=f(x_{2})$ sono detti \emph{minimo} e \emph{massimo} di $f$ in $[a,b]$
(\cref{fig:weierstrass}). Il massimo e il minimo sono unici (\cref{prop:unicita-max}),
i punti di massimo e di minimo in generale no.

\begin{figure}[htb]
\centering
\begin{tikzpicture}
\begin{axis}[width=0.85\textwidth, height=5.2cm, xmin=0, xmax=6.2, ymin=0, ymax=3.4,
  axis lines=left, xtick={0.6,1.62,4.6,5.6}, xticklabels={$a$,$x_{1}$,$x_{2}$,$b$},
  ytick={0.9,3}, yticklabels={$m$,$M$}, xlabel={$x$}, ylabel={$f(x)$}]
  \addplot[curva, domain=0.6:5.6, samples=150, smooth]
    coordinates {(0.6,1.5) (1.0,1.05) (1.62,0.9) (2.4,1.15) (3.2,1.9) (4.0,2.75) (4.6,3.0) (5.2,2.75) (5.6,2.1)};
  \addplot[tratteggio] coordinates {(0,0.9) (1.62,0.9) (1.62,0)};
  \addplot[tratteggio] coordinates {(0,3) (4.6,3) (4.6,0)};
  \addplot[tratteggio] coordinates {(0.6,0) (0.6,1.5)};
  \addplot[tratteggio] coordinates {(5.6,0) (5.6,2.1)};
  \addplot[only marks, mark=*, mark size=1.6pt, rame] coordinates {(1.62,0.9) (4.6,3)};
\end{axis}
\end{tikzpicture}
\caption{Il teorema di Weierstrass: una funzione continua su un intervallo chiuso e
limitato assume un valore minimo $m=f(x_{1})$ e un valore massimo $M=f(x_{2})$.}
\label{fig:weierstrass}
\end{figure}

Il teorema di Weierstrass è dimostrato nel \cref{sec:dim-weierstrass}; in questa sede ci
limitiamo a mettere in luce con degli esempi l'importanza delle ipotesi (funzione
\emph{continua} definita in un intervallo \emph{chiuso} e \emph{limitato}) che
garantiscono l'esistenza del massimo e del minimo (\cref{fig:controesempi-weierstrass}).

\begin{enumerate}[label=\textup{(\alph*)}]
  \item Consideriamo per $x>0$ la funzione $f(x)=\frac1x$. Essa \emph{non assume
  massimo} nell'intervallo aperto a sinistra $(0,1]$; infatti non è limitata
  superiormente in tale intervallo: per ogni $M>0$ risulta $f(x)>M$ se
  $x\in\bigl(0,\frac1M\bigr)$. Qui manca la chiusura dell'intervallo.
  \item La stessa funzione \emph{non assume minimo} nell'intervallo illimitato
  $[1,+\infty)$. La funzione è limitata in tale intervallo, perché $0<\frac1x\le1$ per
  ogni $x\ge1$; però $0$ non è il minimo di $f$ in $[1,+\infty)$, perché non esiste alcun
  $x_{1}\ge1$ per cui $f(x_{1})=\frac1{x_{1}}=0$ (una frazione è nulla se e soltanto se
  il numeratore è nullo), e nessun numero reale positivo $m$ è minimo per $f$, perché
  risulta $f(x)<m$ per ogni $x>\frac1m$. L'estremo inferiore è $0$, ma non è raggiunto:
  qui manca la limitatezza dell'intervallo.
  \item La funzione $f(x)=x^{2}$ assume massimo e minimo in ogni intervallo chiuso e
  limitato $[a,b]$, in particolare nell'intervallo $[-1,1]$. Invece la funzione
  \begin{equation}\label{eq:g-senza-minimo}
    g(x)=
    \begin{cases}
      x^{2} & \text{se } x\in[-1,1]\setminus\{0\},\\
      1 & \text{se } x=0
    \end{cases}
  \end{equation}
  non assume minimo nell'intervallo chiuso e limitato $[-1,1]$, perché assume valori
  positivi arbitrariamente vicini allo zero ($g(x)\to0$ per $x\to0$) ma non è uguale a
  zero per alcun valore di $x\in[-1,1]$. In questo caso la mancanza del minimo è causata
  dalla discontinuità (eliminabile) di $g$ nel punto $x_{0}=0$.
  \item Per la sua discontinuità in corrispondenza dei numeri interi, la funzione
  \emph{parte frazionaria} $f(x)=x-\lfloor x\rfloor$, $x\in\R$, non ha massimo in un
  qualunque intervallo $[a,b]$ che contenga un numero intero $k$ con $a<k\le b$: nei
  punti immediatamente a sinistra di $k$ essa assume valori arbitrariamente vicini a $1$
  (si ha $f(x)=x-k+1$ per $k-1\le x<k$), ma il valore $1$ non è mai assunto, essendo
  $0\le f(x)<1$ per ogni $x$. Si osservi invece che in $[0,\frac12]$, che contiene l'intero
  $0$ ma solo come estremo sinistro, $f(x)=x$ ha massimo $\frac12$.
\end{enumerate}

\begin{figure}[htb]
\centering
\includegraphics[width=0.95\textwidth]{controesempi_weierstrass}
\caption{Quattro funzioni per cui la tesi del teorema di Weierstrass non vale, perché
manca una delle ipotesi: (a) l'intervallo non è chiuso; (b) l'intervallo non è limitato;
(c) e (d) la funzione non è continua (\texttt{python/fig\_controesempi\_weierstrass.py}).}
\label{fig:controesempi-weierstrass}
\end{figure}

Le ipotesi del teorema sono sufficienti, non necessarie: la funzione $\sen x$ ha massimo
$1$ e minimo $-1$ su tutto $\R$, che non è limitato; la funzione di Dirichlet ha massimo
e minimo su ogni intervallo pur essendo discontinua ovunque. Il teorema di Weierstrass
garantisce l'esistenza del massimo, ma non dice dove si trovi: per trovarlo occorrono
altri strumenti (il calcolo differenziale, oppure, come nell'esempio che segue,
un'osservazione algebrica, oppure una ricerca numerica).

\begin{idea}[L'efficienza massima di un velivolo]
Con la \emph{polare parabolica} $C_{D}=C_{D0}+kC_{L}^{2}$, dove $k=\frac1{\pi e\,AR}$
dipende dall'allungamento alare $AR$ e dal fattore di Oswald $e$, l'efficienza
aerodinamica
\[
  E(C_{L})=\frac{C_{L}}{C_{D}}=\frac{C_{L}}{C_{D0}+kC_{L}^{2}}, \qquad 0\le C_{L}\le C_{L,\max},
\]
è una funzione continua su un intervallo chiuso e limitato (il denominatore è positivo):
per il teorema di Weierstrass \emph{esiste} un coefficiente di portanza $C_{L}^{*}$ di
efficienza massima. Per trovarlo, osserviamo che per $C_{L}>0$
\[
  \frac1{E}=\frac{C_{D0}}{C_{L}}+kC_{L}\ \ge\ 2\sqrt{\frac{C_{D0}}{C_{L}}\cdot kC_{L}}=2\sqrt{kC_{D0}},
\]
per la disuguaglianza $u+v\ge2\sqrt{uv}$, valida per $u,v>0$ perché equivale a
$\bigl(\sqrt u-\sqrt v\bigr)^{2}\ge0$, con l'uguaglianza se e solo se $u=v$. Dunque
$E\le E_{\max}=\frac1{2\sqrt{kC_{D0}}}$, con l'uguaglianza se e solo se
$\frac{C_{D0}}{C_{L}}=kC_{L}$, cioè per $C_{L}^{*}=\sqrt{C_{D0}/k}$, la condizione in cui
la resistenza parassita uguaglia quella indotta (purché $C_{L}^{*}\le C_{L,\max}$).
Il programma che segue confronta la formula con una ricerca del massimo su una griglia
fitta, per un velivolo con $C_{D0}=0{,}025$, $AR=8$, $e=0{,}8$ (\cref{fig:efficienza}).
\end{idea}

\begin{python}[Efficienza massima: Weierstrass, poi la formula]
import math

# Polare parabolica: C_D = C_D0 + k C_L^2, con k = 1/(pi e AR)
CD0, AR, e_osw, CLmax = 0.025, 8.0, 0.80, 1.4
k = 1 / (math.pi * e_osw * AR)
E = lambda CL: CL / (CD0 + k * CL**2)       # efficienza L/D

# Weierstrass: E è continua su [0, CLmax], quindi ha massimo.
# Lo cerchiamo su una griglia fitta e lo confrontiamo con la formula.
N = 140000
CL_griglia = max((CLmax * i / N for i in range(N + 1)), key=E)
CL_star = math.sqrt(CD0 / k)
print(f"k = {k:.5f}")
print(f"griglia:  C_L = {CL_griglia:.4f}   E = {E(CL_griglia):.4f}")
E_max = 1 / (2 * math.sqrt(CD0 * k))
print(f"formula:  C_L = {CL_star:.4f}   E = {E_max:.4f}")
\end{python}
\begin{risultato}
k = 0.04974
griglia:  C_L = 0.7090   E = 14.1796
formula:  C_L = 0.7090   E = 14.1796
\end{risultato}

\begin{figure}[htb]
\centering
\includegraphics[width=0.92\textwidth]{efficienza}
\caption{L'efficienza $E=C_{L}/C_{D}$ in funzione di $C_{L}$ per la polare parabolica
dell'esempio: il massimo, che esiste per il teorema di Weierstrass, è
$E_{\max}\approx14{,}18$ per $C_{L}^{*}\approx0{,}709$ (\texttt{python/fig\_efficienza.py}).}
\label{fig:efficienza}
\end{figure}

\subsection{Il secondo teorema dei valori intermedi}

Siamo ora in grado di provare una nuova formulazione del teorema di esistenza dei valori
intermedi.

\begin{teorema}[Secondo teorema dei valori intermedi]\label{teo:valori-intermedi-2}
Una funzione continua in un intervallo $[a,b]$ assume tutti i valori compresi tra il
minimo e il massimo.
\end{teorema}

\begin{schema}
\Ipotesi{$f$ continua in $[a,b]$; \ $m=\min f$, $M=\max f$ (che esistono per Weierstrass).}
\Tesi{per ogni $y_{0}\in[m,M]$ \ $\esiste x_{0}\in[a,b]$: \ $f(x_{0})=y_{0}$.}
\end{schema}

\begin{strategia}
Si ripete la dimostrazione del primo teorema dei valori intermedi, ma al posto degli
estremi $a$, $b$ si usano un punto di minimo $x_{1}$ e un punto di massimo $x_{2}$, che
esistono per il teorema di Weierstrass: la funzione $f(x)-y_{0}$ cambia segno tra
$x_{1}$ e $x_{2}$, e il teorema degli zeri, applicato sull'intervallo di estremi $x_{1}$
e $x_{2}$, fornisce il punto cercato.
\end{strategia}

\begin{dimostrazione}
I valori di massimo $M$ e di minimo $m$ sono assunti in base al teorema di Weierstrass;
rimane da provare che, qualunque sia $y_{0}\in(m,M)$, esiste $x_{0}\in[a,b]$ tale che
$f(x_{0})=y_{0}$. Indichiamo con $x_{1}$, $x_{2}$ i punti di minimo e di massimo di $f$,
cioè tali che $f(x_{1})=m$, $f(x_{2})=M$, e consideriamo la funzione
\begin{equation}\label{eq:vi2-g}
  g(x)=f(x)-y_{0} \qquad\perogni x\in[a,b].
\end{equation}
Essendo $f(x_{1})=m<y_{0}<M=f(x_{2})$, risulta
\begin{equation}\label{eq:vi2-segni}
  g(x_{1})=f(x_{1})-y_{0}<0, \qquad g(x_{2})=f(x_{2})-y_{0}>0;
\end{equation}
per il teorema dell'esistenza degli zeri, applicato a $g$ nell'intervallo di estremi
$x_{1}$, $x_{2}$ (che è $[x_{1},x_{2}]$ oppure $[x_{2},x_{1}]$, e su cui $g$ è continua),
esiste un numero $x_{0}$, appartenente all'intervallo aperto di estremi $x_{1}$, $x_{2}$,
tale che $g(x_{0})=0$, cioè tale che $f(x_{0})=y_{0}$.
\end{dimostrazione}

\begin{corollario}\label{cor:immagine-intervallo}
L'immagine di un intervallo chiuso e limitato mediante una funzione continua è un
intervallo chiuso e limitato: se $f$ è continua in $[a,b]$, allora
$f([a,b])=[m,M]$, dove $m$ ed $M$ sono il minimo e il massimo di $f$.
\end{corollario}

\begin{schema}
\Ipotesi{$f$ continua in $[a,b]$; \ $m=\min f$, $M=\max f$.}
\Tesi{$f([a,b])=[m,M]$.}
\end{schema}

\begin{strategia}
Si dimostrano le due inclusioni: $f([a,b])\subseteq[m,M]$ è la definizione stessa di
minimo e massimo, $[m,M]\subseteq f([a,b])$ è il secondo teorema dei valori intermedi.
\end{strategia}

\begin{dimostrazione}
Ogni valore $f(x)$ sta in $[m,M]$ per definizione di minimo e massimo; viceversa ogni
$y_{0}\in[m,M]$ è un valore di $f$ per il \cref{teo:valori-intermedi-2}.
\end{dimostrazione}

\begin{notastorica}[Una curva che riempie un quadrato]
Una funzione continua trasforma dunque un segmento in un segmento, eventualmente ridotto
a un punto. Che cosa accade se i valori stanno nel piano, cioè se si considerano due
funzioni continue $x(t)$, $y(t)$, con $t\in[0,1]$? L'intuizione suggerisce che il punto
$\bigl(x(t),y(t)\bigr)$ descriva una linea, priva di area. Nel 1878 Georg Cantor aveva
sorpreso i contemporanei dimostrando che tra il segmento e il quadrato esiste una
corrispondenza biunivoca (lo abbiamo raccontato nella nota storica del
\S\ref{sec:invertibili}); l'anno seguente Eugen Netto provò che una tale corrispondenza
non può essere continua. Nel 1890 Giuseppe Peano pubblicò sui \emph{Mathematische
Annalen} una nota di quattro pagine, \emph{Sur une courbe, qui remplit toute une aire
plane}: vi costruiva due funzioni continue $x(t)$, $y(t)$ tali che la coppia
$\bigl(x(t),y(t)\bigr)$, al variare di $t$ in $[0,1]$, assume \emph{tutti} i punti del
quadrato $[0,1]\times[0,1]$. Peano le definì per via aritmetica, con lo sviluppo di $t$
in base tre, senza alcuna figura; l'anno dopo David Hilbert ne pubblicò una variante
geometrica, e fu il primo a illustrarla con un disegno (\cref{fig:curva-hilbert}). Per
il teorema di Netto una curva siffatta non può essere iniettiva: passa più volte per
alcuni punti. La lezione è che la continuità, da sola, non conserva la dimensione; la
conservano le corrispondenze biunivoche e continue insieme alla loro inversa, come
dimostrò Brouwer nel 1911.
\end{notastorica}

\begin{figure}[htb]
\centering
\includegraphics[width=\textwidth]{curva_hilbert}
\caption{Le prime quattro iterate della curva di Hilbert. La $n$-esima è una poligonale
che percorre, uno dopo l'altro, i centri dei $4^{n}$ quadratini di lato $2^{-n}$; ogni
quadratino viene poi diviso in quattro, che l'iterata successiva percorre
consecutivamente. Al crescere di $n$ le poligonali, opportunamente parametrizzate su
$[0,1]$, convergono a una funzione continua che assume tutti i punti del quadrato
(\texttt{python/fig\_curva\_hilbert.py}).}
\label{fig:curva-hilbert}
\end{figure}

\subsection{Un criterio di invertibilità}

Chiudiamo il paragrafo precisando un criterio, introdotto nel \cref{sec:invertibili},
per riconoscere se una data funzione è invertibile. Là avevamo osservato che una
funzione strettamente monotòna è iniettiva (\cref{teo:monotona-invertibile}); la
continuità permette ora di individuarne esattamente l'immagine. La continuità della
funzione inversa è invece studiata nel \cref{sec:monotone-inverse}.

\begin{teorema}[Criterio di invertibilità]\label{teo:invertibilita}\index{funzione!invertibile}
Una funzione continua e strettamente monotòna in un intervallo $[a,b]$ è invertibile in
tale intervallo: se è strettamente crescente, $f:[a,b]\to[f(a),f(b)]$ è biiettiva.
\end{teorema}

\begin{schema}
\Ipotesi{$f$ continua e strettamente crescente (o decrescente) in $[a,b]$.}
\Tesi{$f:[a,b]\to[f(a),f(b)]$ (rispettivamente $[f(b),f(a)]$) è biiettiva.}
\end{schema}

\begin{strategia}
La biiettività si spezza in due: la suriettività sull'intervallo di estremi $f(a)$,
$f(b)$ viene dalla continuità, attraverso il primo teorema dei valori intermedi;
l'iniettività viene dalla stretta monotonia, perché punti distinti hanno immagini
ordinate allo stesso modo, e quindi distinte.
\end{strategia}

\begin{dimostrazione}
Proponiamo la dimostrazione nel caso in cui la funzione $f$ sia strettamente crescente
in $[a,b]$ (il caso decrescente è analogo, con $[f(b),f(a)]$). Risulta
\begin{equation}\label{eq:invertibilita-estremi}
  f(a)<f(x)<f(b) \qquad\perogni x\in(a,b);
\end{equation}
quindi $f(a)$ è il minimo di $f$ in $[a,b]$, mentre $f(b)$ è il massimo, e
$f([a,b])\subseteq[f(a),f(b)]$. Inoltre, per il \cref{teo:valori-intermedi-1}, $f$
assume tutti i valori compresi tra $f(a)$ e $f(b)$: per ogni $y\in[f(a),f(b)]$ esiste
almeno un $x\in[a,b]$ per cui $f(x)=y$. Tale $x$ è unico; infatti, se esistessero due
valori $x_{1}$, $x_{2}$ distinti tra loro, diciamo $x_{1}<x_{2}$, per cui
$y=f(x_{1})=f(x_{2})$, allora dovrebbe risultare anche $f(x_{1})<f(x_{2})$, dato che $f$ è
strettamente crescente. Quindi $f:[a,b]\to[f(a),f(b)]$ è suriettiva e iniettiva, cioè
invertibile.
\end{dimostrazione}
